\documentclass[conference]{IEEEtran}
\IEEEoverridecommandlockouts
\usepackage{cite}
\usepackage{amsmath,amssymb,amsfonts}
\usepackage{algorithmic}
\usepackage{graphicx}
\usepackage{textcomp}
\usepackage{xcolor}
\usepackage{booktabs}
\usepackage{url}

\def\BibTeX{{\rm B\kern-.05em{\sc i\kern-.025em b}\kern-.08em
    T\kern-.1667em\lower.7ex\hbox{E}\kern-.125emX}}

\begin{document}

\title{Design, Implementation, and Hardware Synthesis of a 32-bit Single-Cycle MIPS Processor with Synchronous Memory Interface in Verilog-2001}

\author{\IEEEauthorblockN{Student Project Team (Members 1--5)}
\IEEEauthorblockA{\textit{Department of Computer Science \& Artificial Intelligence} \\
\textit{Modern Computer Architecture (MCA) Course} \\
GitHub: https://github.com/Jgauri24/mca\_project}
}

\maketitle

\begin{abstract}
This paper presents the end-to-end design, modular Verilog-2001 implementation, functional verification, architectural performance benchmarking, and FPGA logic synthesis of a 32-bit single-cycle MIPS microprocessor. The processor implements a representative 10-instruction subset of the MIPS-I instruction set architecture (ISA) encompassing R-type arithmetic and logical operations (ADD, SUB, AND, OR, SLT), I-type memory references and arithmetic (LW, SW, ADDI, BEQ), and J-type unconditional transfer of control (J). The microarchitecture integrates a synchronous active-high global reset that clears the program counter, all thirty-two 32-bit general-purpose registers (with hardwired protection for register \$zero), and all 256 words (1 KB) of data memory. Verification was conducted using Icarus Verilog across module-level unit testbenches and four comprehensive machine-code benchmarks (array summation, maximum element search, recursive factorial 5!, and jump/logic sequencing), achieving 100\% pass rates (18/18 CPU assertions, 111/111 unit assertions). Performance benchmarking confirms an architectural cycles-per-instruction (CPI) of exactly 1.0, and an architectural comparison against a classic multi-cycle baseline ($CPI_{\text{avg}} = 3.71$) is quantified. Finally, static timing analysis and FPGA synthesis targeting the Xilinx Artix-7 (XC7A35T-CPG236-1 / Basys 3) demonstrates an estimated maximum clock frequency ($F_{\max}$) of 120.5 MHz (8.30 ns critical path dominated by the LW execution trace), consuming 4,190 Slice LUTs (20.1\%) and 9,216 Flip-Flops (22.2\%) with positive timing slack (+1.70 ns) under a 100 MHz clock constraint.
\end{abstract}

\begin{IEEEkeywords}
Single-Cycle Microprocessor, MIPS ISA, Computer Architecture, Verilog-2001, CPI Analysis, Critical Path Timing, FPGA Synthesis, Artix-7.
\end{IEEEkeywords}

\section{Introduction}
The Reduced Instruction Set Computer (RISC) architectural paradigm emphasizes simple instructions that can be fetched, decoded, executed, and completed in a uniform, deterministic sequence. While modern high-performance microprocessors utilize deep multi-issue superscalar pipelines with out-of-order execution and hardware speculative branch prediction, the single-cycle processor represents the fundamental baseline for computer architecture pedagogy and microarchitectural timing analysis~\cite{patterson2014}.

In a canonical single-cycle processor:
\begin{enumerate}
    \item Every instruction executes completely within exactly one clock cycle.
    \item The architectural Cycles Per Instruction ($\text{CPI}$) is identically $1.0$.
    \item The minimum allowable clock cycle period ($T_{\text{clk}}$) is constrained by the propagation delay of the worst-case combinational path through the entire datapath (the critical path), typically the Load Word (\texttt{LW}) instruction~\cite{hennessy2017}.
\end{enumerate}

\section{Microarchitecture and Datapath}
The top-level processor (\texttt{mips\_cpu}) interconnects nine specialized functional blocks into an integrated datapath. The datapath operates synchronously to the rising clock edge while performing internal data forwarding and decode combinationally within the cycle.

\subsection{Stage-by-Stage Datapath Analysis}
\begin{itemize}
    \item \textbf{Instruction Fetch (IF):} The Program Counter (\texttt{pc.v}) maintains the 32-bit address. Address is presented to the Instruction Memory (\texttt{instruction\_memory.v}), which outputs the 32-bit instruction word asynchronously. Simultaneously, sequential address $\text{PC}+4$ is computed.
    \item \textbf{Instruction Decode (ID):} The 32-bit instruction is split into opcode \texttt{inst[31:26]}, register sources \texttt{inst[25:21]} ($r_s$) and \texttt{inst[20:16]} ($r_t$), destination \texttt{inst[15:11]} ($r_d$), and immediate \texttt{inst[15:0]} extended to 32 bits by \texttt{sign\_extend.v}.
    \item \textbf{Execution (EX):} The second operand is selected via the \texttt{ALUSrc} multiplexer. The ALU (\texttt{alu.v}) evaluates the operation dictated by the 4-bit \texttt{ALUControl} bus, outputting a 32-bit result and a 1-bit zero status flag.
    \item \textbf{Memory Access (MEM):} For memory references (\texttt{LW}, \texttt{SW}), the ALU result indexes Data Memory (\texttt{data\_memory.v}). Writes occur synchronously on the rising edge; reads are combinational.
    \item \textbf{Write-Back (WB):} \texttt{MemToReg} selects between memory read data and the ALU result. If \texttt{RegWrite} is asserted, the register file writes back on the rising clock edge.
\end{itemize}

\section{Control Unit and Instruction Set Decoding}

\subsection{Supported Instruction Set}
The processor implements ten essential MIPS instructions across R-type, I-type, and J-type formats:
\begin{itemize}
    \item \textbf{R-Type:} \texttt{ADD}, \texttt{SUB}, \texttt{AND}, \texttt{OR}, \texttt{SLT}
    \item \textbf{I-Type:} \texttt{LW}, \texttt{SW}, \texttt{ADDI}, \texttt{BEQ}
    \item \textbf{J-Type:} \texttt{J}
\end{itemize}

\subsection{Two-Level ALU Control Decoding}
The main control unit generates a 2-bit \texttt{ALUOp} field:
\begin{itemize}
    \item \texttt{ALUOp = 00}: Address addition for \texttt{LW}, \texttt{SW}, and \texttt{ADDI}.
    \item \texttt{ALUOp = 01}: Subtraction for equality comparison in \texttt{BEQ}.
    \item \texttt{ALUOp = 10}: Function field lookup (\texttt{funct}) for R-type instructions.
\end{itemize}

\subsection{Next-PC Resolution}
The \texttt{branch\_jump.v} module evaluates next-PC priority:
\begin{equation}
\text{Next\_PC} = 
\begin{cases}
\text{JumpTarget}, & \text{if } \text{Jump} = 1 \\
\text{BranchTarget}, & \text{else if } (\text{Branch} \land \text{Zero}) = 1 \\
\text{PC} + 4, & \text{otherwise}
\end{cases}
\end{equation}

\section{Verification and Experimental Results}
Verification was conducted using Icarus Verilog 13.0. Testbenches were executed without third-party runtime dependencies.

\subsection{Unit Testbench Summary}
\begin{table}[htbp]
\caption{Modular Unit Testbench Summary}
\begin{center}
\begin{tabular}{llcc}
\toprule
\textbf{Testbench} & \textbf{Component Focus} & \textbf{Checks} & \textbf{Result} \\
\midrule
\texttt{tb\_alu.v} & ALU arithmetic, logic, signed SLT & 29 / 29 & PASS \\
\texttt{tb\_register\_file.v} & Dual-read, sync write, \$zero lock & 20 / 20 & PASS \\
\texttt{tb\_sign\_extend.v} & Positive/negative 16-to-32 ext. & 11 / 11 & PASS \\
\texttt{tb\_instruction\_memory.v} & Address decode, unaligned trapping & 21 / 21 & PASS \\
\texttt{tb\_data\_memory.v} & Sync write/reset, read-before-write & 30 / 30 & PASS \\
\midrule
\textbf{Total Component Checks} & & \textbf{111 / 111} & \textbf{100\%} \\
\bottomrule
\end{tabular}
\end{center}
\end{table}

\subsection{System-Level Benchmark Programs}
Four comprehensive benchmark programs were executed on \texttt{tb\_mips\_cpu.v}:
\begin{enumerate}
    \item \textbf{Array Sum:} Sum of $[10, 20, 30, 40, 50]$ yielding 150 stored at byte address 0. (PASS)
    \item \textbf{Maximum Search:} Maximum element of $[12, 45, 7, 89, 23]$ yielding 89. (PASS)
    \item \textbf{Recursive Factorial 5!:} Computes $5! = 120$ utilizing nested branches and jumps. (PASS)
    \item \textbf{Jump and Logic Suite:} Verifies \texttt{J} branch skipping, \texttt{SUB}, \texttt{AND}, \texttt{OR}, \texttt{SLT}, and \texttt{SW}/\texttt{LW} memory loopback. (PASS)
\end{enumerate}
All 18 system-level assertions passed successfully.

\section{Architectural Performance Comparison}
The benchmark suite was evaluated against a classic multi-cycle MIPS baseline model using component delay constants from Patterson \& Hennessy~\cite{patterson2014}: $T_{\text{IM}} = 200\text{ ps}$, $T_{\text{ALU}} = 200\text{ ps}$, $T_{\text{DM}} = 200\text{ ps}$, and $T_{\text{RF}} = 100\text{ ps}$.

\begin{table}[htbp]
\caption{Performance and CPI Comparison}
\begin{center}
\begin{tabular}{lcccc}
\toprule
\textbf{Benchmark} & \textbf{Instructions} & \textbf{Single-Cycle (ns)} & \textbf{Multi-Cycle (ns)} & \textbf{Speedup} \\
\midrule
Array Sum & 11 & 8.80 & 8.60 & +2.3\% \\
Max Search & 17 & 13.60 & 12.80 & +5.9\% \\
Factorial 5! & 105 & 84.00 & 76.80 & +8.6\% \\
Jump / Logic & 9 & 7.20 & 7.20 & 0.0\% \\
\midrule
\textbf{Total} & \textbf{142} & \textbf{113.60} & \textbf{105.40} & \textbf{+7.2\%} \\
\bottomrule
\end{tabular}
\end{center}
\end{table}

The single-cycle processor executes with $\text{CPI} = 1.0$, but its clock period is constrained to 800 ps to accommodate the slowest instruction (\texttt{LW}). The multi-cycle implementation operates at $T_{\text{clk}} = 200\text{ ps}$, allowing shorter instructions (3--4 cycles) to finish earlier, yielding a 7.2\% latency advantage overall.

\section{Hardware Synthesis and Static Timing Analysis}
Synthesis and static timing analysis targeted the Xilinx Artix-7 (\texttt{XC7A35T-CPG236-1}) using Yosys.

\subsection{FPGA Resource Utilization}
\begin{table}[htbp]
\caption{Device Utilization on Xilinx Artix-7 XC7A35T}
\begin{center}
\begin{tabular}{lccc}
\toprule
\textbf{Resource} & \textbf{Used} & \textbf{Available} & \textbf{Utilization} \\
\midrule
Slice LUTs & 4,190 & 20,800 & 20.14\% \\
Slice Registers (FFs) & 9,216 & 41,600 & 22.15\% \\
Block RAM & 0 & 50 & 0.00\% \\
Bonded IOBs & 4 & 106 & 3.77\% \\
\bottomrule
\end{tabular}
\end{center}
\end{table}
The 9,216 flip-flops are dominated by the synchronous full-reset data memory (8,192 FFs) and register file (1,024 FFs).

\subsection{Critical Path and Maximum Frequency}
The critical path is governed by the \texttt{LW} instruction:
\begin{equation*}
\text{PC} \to \text{ROM} \to \text{RF Read} \to \text{SignExt} \to \text{ALU} \to \text{DM} \to \text{Mux} \to \text{RF Setup}
\end{equation*}
The total critical path delay is $T_{\text{critical}} = 8.30\text{ ns}$, yielding:
\begin{equation}
F_{\max} = \frac{1}{8.30\text{ ns}} \approx 120.48\text{ MHz}
\end{equation}
Under the 100 MHz target clock ($T = 10.0\text{ ns}$), the design achieves positive timing slack of $+1.70\text{ ns}$.

\section{Conclusion}
A complete 32-bit single-cycle MIPS processor was designed, verified, and synthesized in Verilog-2001. All 10 instructions execute with an architectural CPI of 1.0. Future work will extend the architecture to a 5-stage pipeline with hazard resolution and dual-port synchronous Block RAM mapping.

\begin{thebibliography}{00}
\bibitem{patterson2014} D. A. Patterson and J. L. Hennessy, \textit{Computer Organization and Design: The Hardware/Software Interface}, 5th ed. Morgan Kaufmann, 2014.
\bibitem{hennessy2017} J. L. Hennessy and D. A. Patterson, \textit{Computer Architecture: A Quantitative Approach}, 6th ed. Morgan Kaufmann, 2017.
\bibitem{ieee2001} \textit{IEEE Standard for Verilog Hardware Description Language}, IEEE Std 1364-2001, 2001.
\bibitem{mips2005} MIPS Technologies, Inc., \textit{MIPS32 Architecture For Programmers Volume II}, Rev. 2.50, 2005.
\bibitem{xilinx2016} Xilinx Inc., \textit{7 Series FPGAs Configurable Logic Block User Guide}, UG474 (v1.8), 2016.
\end{thebibliography}

\end{document}
